% 第6章第4节：一般方阵的对角化（复特征值与谱半径、实Schur形与正规矩阵）。

\section{一般方阵的对角化}
\label{sec:ls-general-square-spectrum}

前两节的对角化各有前提：用可逆矩阵对角化，要求有足够多的实特征向量；用正交矩阵对角化，要求矩阵对称。不对称的方阵会在两处偏离这些前提。其一，实特征值可能不够：平面上的$90^\circ$旋转没有实特征向量，第~\ref{sec:ls-matrix-powers}节的对角化在第一步就停止。其二，即使可以对角化，特征向量也可能互不正交：示例~\ref{ex:ls-triangular-power}中矩阵的两条特征方向$(1,0)^\top$与$(-40/3,1)^\top$夹角只有约$4.3^\circ$（计算见示例~\ref{ex:ls-nonnormal-norm-radius}）。本节回答三个问题：非实的特征值在实空间中对应什么作用，怎样用一个数概括矩阵每作用一次的伸缩倍数；若坚持用正交矩阵换基，一般方阵最多能化到多接近对角的形式；什么样的方阵在这种换基下仍能化为分块对角形。

\subsection{复特征值与谱半径}

第~\ref{sec:ls-characteristic-polynomial}节的特征多项式$p_{\bm A}$是$d$次实系数多项式。由代数基本定理，它在复数范围内恰有$d$个根（计重数）。把这些根都称为$\bm A$的特征值，全体特征值组成的集合称为$\bm A$的\term{谱}（spectrum）。由于系数都是实数，若$p_{\bm A}(\lambda)=0$，两边取共轭即得$p_{\bm A}(\bar\lambda)=0$，所以非实根成共轭对出现。实特征值的情形与第~\ref{sec:ls-linear-operators}节相同。

第~\ref{sec:ls-eigen-similarity-invariance}节由系数比较得到$\operatorname{tr}\bm A=\sum_i\lambda_i$与$\det\bm A=\prod_i\lambda_i$，那里要求$p_{\bm A}$能分解为实一次因子。计入复特征值以后，$p_{\bm A}$在复数范围内总能分解为一次因子之积，同样的系数比较说明，这两个等式对每个方阵都成立。

非实特征值在实空间里对应什么作用？设$\lambda=a+\mathrm ib$（$b\neq0$）是特征值。复矩阵$\bm A-\lambda\bm I$的行列式为零，所以有非零的复向量$\bm u+\mathrm i\bm w$满足$\bm A(\bm u+\mathrm i\bm w)=(a+\mathrm ib)(\bm u+\mathrm i\bm w)$，其中$\bm u,\bm w$为实向量。比较两边的实部与虚部，得
\[
 \bm A\bm u=a\bm u-b\bm w,\qquad \bm A\bm w=b\bm u+a\bm w .
\]

这两个等式说明$\operatorname{span}\{\bm u,\bm w\}$在$\bm A$下不变。它是二维的，即$\bm u,\bm w$线性无关。否则$\bm u,\bm w$（不全为零）只张成一条直线$\operatorname{span}\{\bm y\}$，这条直线在$\bm A$下不变，所以$\bm A\bm y=\mu\bm y$，$\mu$为实数；而$\bm u+\mathrm i\bm w$是$\bm y$的复数倍，于是$\bm A(\bm u+\mathrm i\bm w)=\mu(\bm u+\mathrm i\bm w)$，得$\lambda=\mu$为实数，与$b\neq0$矛盾。所以$\operatorname{span}\{\bm u,\bm w\}$是二维实不变子空间。在基$(\bm u,\bm w)$下，$\bm A$的作用由矩阵
\[
 \begin{bmatrix}a&b\\-b&a\end{bmatrix}
 =|\lambda|\begin{bmatrix}\cos\phi&\sin\phi\\-\sin\phi&\cos\phi\end{bmatrix},
 \qquad \cos\phi=\frac{a}{|\lambda|},\ \sin\phi=\frac{b}{|\lambda|}
\]
表示，即$|\lambda|$倍的伸缩加一次旋转。

所以无论特征值是实是复，它的模长$|\lambda|$都是沿相应方向或平面每作用一次的伸缩倍数；复特征值的情形在基$(\bm u,\bm w)$的坐标中度量。全部特征值模长的最大值
\[
 \rho(\bm A)=\max_i|\lambda_i(\bm A)|
\]
称为\term{谱半径}（spectral radius）。例如$90^\circ$旋转的特征值为$\pm\mathrm i$，谱半径为$1$：它不改变长度，只改变方向。

\subsection{实Schur形与正规矩阵}

\paragraph{实Schur形}

一般方阵未必能对角化，更未必能用正交矩阵对角化。退一步，能否用正交矩阵把它化为上三角形？上三角形同样便于计算：对角元就是特征值，解方程和求幂都可以从最后一个坐标逐个向前推。实特征值不够时，实上三角形不可能存在，因为它的对角元就是特征值，必须全是实数；但允许对角线上出现$2\times2$块以后，答案总是肯定的。

\begin{theorem}[实Schur形]
\label{thm:ls-real-schur}
对任意$\bm A\in\mathbb R^{d\times d}$，存在正交矩阵$\bm Q$与分块上三角矩阵$\bm T$，使
\begin{equation}
 \bm A=\bm Q\bm T\bm Q^{\top},
 \label{eq:matrix-analysis-schur-decomposition}
\end{equation}
$\bm T$的对角块为$1\times1$（实特征值）或$2\times2$（一对共轭特征值），对角块下方全为零。$\bm A$的特征值就是各对角块的特征值。
\end{theorem}
\begin{proof}
对$d$归纳，思路与谱定理相同：先找一个不变子空间，再对其正交补降维。若$\bm A$有实特征值$\lambda$，取单位特征向量$\bm q_1$并扩充为正交矩阵$\bm Q_1=[\bm q_1,\bm Q_2]$。由$\bm A\bm q_1=\lambda\bm q_1$与$\bm Q_2^\top\bm q_1=\bm0$，$\bm Q_1^\top\bm A\bm Q_1$的第一列为$\lambda\bm e_1$，即
\[
 \bm Q_1^\top\bm A\bm Q_1=\begin{bmatrix}\lambda&*\\\bm0&\bm Q_2^\top\bm A\bm Q_2\end{bmatrix}.
\]
若没有实特征值，取上一小节得到的二维不变子空间，用它的一组标准正交基$\bm q_1,\bm q_2$同样扩充，得到左上角为$2\times2$块、其下方为零的分块形式，该块的特征值是$\lambda,\bar\lambda$。分块上三角矩阵的特征多项式是对角块特征多项式之积，所以右下块的特征值恰为其余特征值。对右下块应用归纳假设，再把两次正交变换合并，即得结论。
\end{proof}
实际计算先用Householder反射化为上Hessenberg形，再作带位移的QR迭代，标准数值软件都提供这一分解；本卷只使用它的存在性。第~\ref{sec:ls-sylvester-schur}节求解Sylvester方程时，用的正是这一形式。

\paragraph{正规矩阵}

Schur形什么时候是分块对角的？实对称矩阵是一例。这时$\bm T=\bm Q^\top\bm A\bm Q$仍对称，它的$2\times2$对角块也对称，特征值为实数，不可能对应一对共轭特征值，所以$\bm T$只有$1\times1$对角块，是上三角矩阵；对称的上三角矩阵只能是对角矩阵，这又一次得到谱定理。更一般地，第~\ref{sec:ls-adjoints}节把满足$TT^\star=T^\star T$的算子称为正规算子。在标准坐标下伴随就是转置，满足$\bm A^\top\bm A=\bm A\bm A^\top$的矩阵称为\term{正规矩阵}（normal matrix）；对称、反对称与正交矩阵都是正规的。

正规矩阵的Schur形$\bm T$分块对角。首先$\bm T$也正规：$\bm T^\top\bm T=\bm Q^\top\bm A^\top\bm A\bm Q=\bm Q^\top\bm A\bm A^\top\bm Q=\bm T\bm T^\top$。比较$\bm T^\top\bm T$与$\bm T\bm T^\top$的左上角对角块：前者是$\bm T_{1,1}^\top\bm T_{1,1}$，因为$\bm T$的第一块列只有$\bm T_{1,1}$非零；后者是$\sum_j\bm T_{1,j}\bm T_{1,j}^\top$。两边取迹，得
\[
 \lVert\bm T_{1,1}\rVert_{\mathrm F}^2=\lVert\bm T_{1,1}\rVert_{\mathrm F}^2+\sum_{j>1}\lVert\bm T_{1,j}\rVert_{\mathrm F}^2 ,
\]
故第一块行的非对角块全为零。去掉第一块后，剩下的矩阵仍正规，逐块推下去可知$\bm T$分块对角。

每个$2\times2$对角块$\left[\begin{smallmatrix}p&q\\r&s\end{smallmatrix}\right]$也正规。比较两边的$(1,1)$元得$q^2=r^2$。若$q=r$，块对称，特征值为实数，与$2\times2$块对应一对共轭特征值矛盾；所以$r=-q\neq0$。再比较$(1,2)$元得$2q(s-p)=0$，即$s=p$。于是每块都形如$\left[\begin{smallmatrix}a&b\\-b&a\end{smallmatrix}\right]$，是模长乘旋转。反过来，若正交矩阵$\bm Q$把$\bm A$化为分块对角矩阵$\bm T$，且每块是一个数或这种模长乘旋转的块，则每块与自己的转置可交换，$\bm T^\top\bm T=\bm T\bm T^\top$，从而$\bm A=\bm Q\bm T\bm Q^\top$正规。所以正规矩阵恰好是能用标准正交基化为这种分块对角形的实方阵。

总之，正规矩阵沿一组标准正交的方向或平面只作伸缩，或伸缩加旋转。$\bm T^k$仍分块对角，每块是$\lambda^k$或$|\lambda|^k$乘一个旋转，二范数为$|\lambda|^k$；正交变换又不改变谱范数，所以$\lVert\bm A^k\rVert_2=\lVert\bm T^k\rVert_2=\max_i|\lambda_i|^k=\rho(\bm A)^k$。也就是说，正规矩阵每一步的最大放大倍数恰好就是谱半径。

\paragraph{非正规矩阵}
由上一段，非正规矩阵的Schur形$\bm T$不可能是由数与模长乘旋转块组成的分块对角形：它要么含非零的非对角块，要么某个$2\times2$对角块本身不正规。非正规矩阵也不能用一组标准正交的特征向量对角化，否则它对称。这时矩阵范数可以远大于谱半径，下面的例子显示了这种分离。
\begin{example}[非正规矩阵的范数与谱半径]
\label{ex:ls-nonnormal-norm-radius}
给定示例~\ref{ex:ls-triangular-power}的$\bm A=\left[\begin{smallmatrix}0.8&4\\0&0.5\end{smallmatrix}\right]$，比较它的谱半径、一步的放大倍数与特征方向的夹角。

它不正规：$\bm A^\top\bm A$的$(1,1)$元为$0.8^2=0.64$，$\bm A\bm A^\top$的$(1,1)$元为$0.8^2+4^2=16.64$。特征值为$0.8$与$0.5$，所以$\rho(\bm A)=0.8$。一步作用把$(0,1)^\top$送到$(4,0.5)^\top$，长度为$\sqrt{16.25}>4$，所以$\lVert\bm A\rVert_2>4$。两条特征方向$\bm v_1=(1,0)^\top$与$\bm v_2=(-40/3,1)^\top$的夹角$\theta$满足
\[
 \cos\theta=\frac{|\bm v_1^\top\bm v_2|}{\lVert\bm v_1\rVert_2\lVert\bm v_2\rVert_2}
 =\frac{40/3}{\sqrt{(40/3)^2+1}}\approx0.997,
\]
即$\theta\approx4.3^\circ$。结论：两条特征方向几乎平行，一步的最大放大倍数超过$4$，而谱半径只有$0.8$。
\end{example}
谱半径描述反复作用的长期倍率，范数描述一步的最大放大；非正规时二者分离。第~\ref{sec:ls-matrix-power-behavior}节将看到，这会造成先放大、后衰减的瞬态。

至此，本节的问题有了回答。实矩阵的特征值可能是复数，一对共轭特征值$a\pm\mathrm ib$对应一个二维不变平面上的伸缩加旋转；谱半径$\rho(\bm A)$概括每作用一次的长期伸缩倍数。坚持用标准正交基时，一般方阵只能化为实Schur形；恰好能化为由数与模长乘旋转块组成的分块对角形的是正规矩阵，它每一步的最大放大等于谱半径，非正规矩阵的范数与谱半径则可能相差悬殊。这些结论都要求矩阵把空间映到自身，并且输入、输出用同一组基。若允许两端各取一组标准正交基，这一限制就消失了：下一节说明，任何矩阵都能这样化为对角形。
